% If you want to compile this section only, make sure to include relevant document headers and the \being \end document commands.
% You can make this a bit easier if you use the subfile package

Misinformation on social media is difficult because it is fast. Posts can reach large audiences before strong evidence about accuracy is available, and by the time a claim is confidently verified or debunked, the main burst of attention may have already passed. This creates a timing problem: decisions about whether to flag a claim or prioritize it for verification often have to be made while the cascade is still unfolding, using only a partial snapshot of early spread \citep{colbaugh_earlywarning_2012}. A COVID-era misinformation dataset illustrates the stakes. Two false claims can look relatively quiet in the first hour—showing only shallow early diffusion—yet end up with dramatically different reach. Despite their massive difference in final reach, their first-hour diffusion looks comparable in scale, making it hard to distinguish likely high-reach cascades from early signals alone. One claim—“trump bailed out wall street but not main street”—eventually grows into a massive cascade (final size $=15096$), while another—“church services cant resume until were all vaccinated says bill gates”—quickly stalls (final size $=16$). This contrast motivates explicit early-observation windows: the practical question is not whether we can judge truth perfectly in hindsight, but whether we can identify potentially high-risk cases early enough to prioritize limited attention and verification resources. Beyond operational triage, this is also a social-science question: early attention and reinforcement dynamics may shape a cascade's trajectory before reliable evidence about accuracy becomes available. 


Most existing misinformation detection work focuses on content- and user-based features \citep{shu_fake_2017,zubiaga_detection_2018}. These approaches can be powerful once a post has accumulated enough text context, reactions, and external signals, but early on those signals are often sparse or missing, making predictions noisier---precisely when rapid spread leaves little time for verification. This project explores a complementary source of information: the early diffusion process itself. If sharing behavior is shaped by social influence and network constraints, then the first part of a cascade may already contain predictive signals about subsequent spread \citep{szabo_predicting_2010,cheng_can_2014}; whether those early signals also help with veracity prediction is an empirical question.

To examine this question, I study early diffusion patterns as signals for two related outcomes. The first is a diffusion-outcome task: can we predict how far a cascade will spread (e.g., final size or depth) from its initial trajectory? The second is a credibility task: can we use early diffusion information to identify cascades that are at higher misinformation risk? There is evidence that true and false information diffuse differently on social platforms \citep{vosoughi_spread_2018}, but much of that evidence is retrospective, based on fully observed cascades \citep{del_vicario_spreading_2016}. In the next section, I formalize this early-warning setting with explicit early-observation windows and evaluate how performance changes as those windows expand.


Accordingly, the central research question is not only whether early diffusion can predict outcomes, but \emph{when} prediction becomes reliably informative under explicit early-observation windows. Throughout, I treat size-based windows (the first $k$ messages) as the primary setting and report time-based windows as robustness checks reflecting real-time constraints. A secondary question is whether structural and temporal diffusion features add signal beyond simple early-activity baselines within the same window.

Operationally, I define “early” in two complementary ways: a time-based window (e.g., the first 60 minutes after the source post) and a size-based window (e.g., the first $k$ messages). Using both allows me to test whether early diffusion patterns add information beyond simple early activity and whether those gains differ across the reach and veracity tasks.

This project sits at the intersection of two strands of work. One line studies early prediction of cascade outcomes and shows that forecasting eventual reach from partial trajectories is possible but difficult \citep{cheng_can_2014,szabo_predicting_2010}. A second line documents systematic diffusion differences between true and false information, largely using fully observed cascades \citep{vosoughi_spread_2018}. What remains less clear is how early these differences become actionable under a realistic early-warning constraint---that is, when models are restricted to the first $T$ minutes (e.g., 30/60/180 minutes) or the first $k$ messages of a cascade. This connects these strands by testing whether diffusion-based signals provide predictive value within such early windows, and whether the timing of informativeness differs between reach prediction and veracity discrimination.

Empirically, the project uses propagation datasets with timestamps, parent--child links, and veracity labels. FibVID serves as the main dataset, while Twitter15/16 are retained as a comparative benchmark because they provide the original native rumor-tree setting for the early-diffusion pipeline. From each early window, I construct structural and temporal features that summarize partial propagation, including early volume (tweets/users), partial depth and width, branching patterns, and early growth dynamics based on interarrival times.

The main empirical pattern that emerges is a strong asymmetry between the two tasks: early diffusion is much more informative for eventual reach than for veracity, while veracity results are better interpreted as weak screening evidence than as strong classification performance.

The paper proceeds as follows. Section~2 reviews work on early-stage diffusion dynamics, mechanisms of social influence, and the gap between retrospective diffusion differences and early-warning prediction. Section~3 describes the datasets, early-window definitions, feature construction, and evaluation plan. Sections~4 and~5 present results and discuss implications and limitations for early warning based on diffusion signals.
